% Réplication triple contre codage d'effacement Reed-Solomon (10, 4).
\documentclass[border=6pt]{standalone}
\usepackage{sd}
\begin{document}
\begin{tikzpicture}[x=1mm, y=1mm, note/.append style={font=\sffamily\small},
    f/.style={draw=#1, fill=#1Bg, line width=0.5pt, minimum width=7mm, minimum height=7mm, font=\sffamily\tiny, inner sep=0pt}]
  \node[pillo=Dblue] at (40,56) {réplication triple};
  \foreach \k in {0,1,2} \node[f=Dblue, minimum width=22mm, font=\sffamily\footnotesize] at ({4+\k*27},44) {copie \the\numexpr\k+1\relax};
  \node[note, anchor=north, align=center] at (31,37) {3 fois la place ; supporte la perte de 2 copies ;\\une lecture lit une copie, une réparation recopie tout l'objet};
  \begin{scope}[yshift=-38mm]
    \node[pillo=Dgreen] at (40,56) {codage d'effacement Reed-Solomon (10, 4)};
    \foreach \k in {0,...,9} \node[f=Dgreen] at ({\k*7.5},44) {D\k};
    \foreach \k in {0,...,3} \node[f=Dorange] at ({78+\k*7.5},44) {P\k};
    \draw[decorate, decoration={brace, amplitude=2mm, mirror}, line width=0.6pt] (-3.5,39) -- (71,39);
    \node[note, anchor=north] at (34,36) {10 fragments de données};
    \draw[decorate, decoration={brace, amplitude=2mm, mirror}, line width=0.6pt] (74.5,39) -- (104,39);
    \node[note, anchor=north] at (89,36) {4 de parité};
    \node[note, anchor=north, align=center] at (50,28) {1,4 fois la place ; n'importe quels 10 fragments sur 14 suffisent ;\\une réparation lit 10 fragments pour en recalculer un};
  \end{scope}
\end{tikzpicture}
\end{document}
